\section{Introduction}

In type theory, universes allow manipulating types as terms.
To avoid Girard's paradox, it is routine to introduce a \textit{hierarchy} of universes: $\U_0, \U_1, \U_2, ...$.
In practical systems, those are often assumed to be \textit{cumulative}, requiring proper inclusions between them.
There are two different ways to formulate the notion of universe in type theory. Martin-Löf’s
first versions \cite{MARTINLOF75, MARTINLOF82} were formulated à la Russell, where a universe is seen as a type of types. Later, an alternative version à la Tarski \cite{MARTINLOF84}, where a universe is seen as a type of codes along with a decoding function, was also introduced. Usually, Russell universes are regarded as an informal notation of the latter, as \cite{Luo} suggested.

When presenting Type Theory in an algebraic fashion, Tarski universes are essential, at least if we do not want to introduce sort equations, and one can show that the system à la Tarski presents the initial model of such a theory.
This initial model has been constructed explicitly in \citet{bcde2026}, following the method presented in \citet{Streicher} and formalised in \citet{Brunerie}. 
However, there are as many initiality theorems as there are presentations of type theory, and if people often assume that they
extend easily to the more informal syntax used in practice, this comes from misconceptions as they depend on subtle details of the presentation.
To solve this issue, Voevodsky proposed the \textit{Initiality Conjecture project}, which aimed at developing rigorous metatheory for large class of theories.
A first step at solving this goal was achieved by Streicher \cite{Streicher}, proving initiality for a fully annotated Calculus of Constructions, and then
proving its equivalence with the more common, unannotated formulation, close to a Pure Type System. 

Generalised Algebraic Theories (GATs) provide a framework to describe type theories and their model in a uniform
way, and to prove initiality rather easily \cite{noteGAT, Brunerie}, 
thus solving the first part of the problem. However, GATs give rise to a fully explicit and annotated presentation of type theory, which is inconvenient to work in.
Therefore, it would be essential to extend Streicher's equivalence results to relate Generalised Algebraic and Pure Type System presentations, where the latter is presented
with Russell universes and unannotated syntax.  This would justify rigorously the common practice of using informal syntax and Russell universes, as proposed by \citet{Luo},
and would legitimate interpreting such syntax in the usual models of type theory.
In particular, our work essentially relates the results proven on explicit type theories using semantic arguments, to the corresponding implicit theories implemented in proof systems.

\noindent
\textbf{Related work.} In recent work, Assaf and Thiré \cite{assaf,thire} show the equivalence between Russell and Tarski presentation of universes for the Calculus
of Constructions with $\beta$ as an external reduction. Nevertheless, it appeared that their proof was 
incorrect, and a new proof using the same technique was finally published by \citet{tiago-fscd}. Unfortunately, their arguments do not apply to
formulations of type theory with judgemental equality, and do not extend to $\eta$-equivalence. 
This legitimates the study of a new proof technique. 

In the present paper, we will not consider cumulativity with a structural sub-typing judgement $A \leq B$ \cite{assaf,thire,Luo, subtyping},
as it can lead to discrepancy between Russell and Tarski universes in a more general setting, as pointed out by Luo \cite{Luo}.
For instance, sub-typing breaks canonicity when adding inductive types, and does not hold for $\Pi$-types in the Set model, or it is at least unclear
depending on the encoding of application in the ambient set theory. 
Therefore, here we will study the simplest form of cumulativity that is possible.
Yet, structural sub-typing would provide a more convenient manipulation of universes in a practical proof
assistant. Notably, it simplifies greatly the use of algebraic hierarchies \cite{subtyping}.
The recent work of \citet{AdapTT} presents a general framework which would provide automatically the necessary
equations to solve the issues mentioned above. In any case, we believe that our work is orthogonal to such considerations,
and that our proof should extend to any system that handles structural sub-typing in a correct way.\\

\noindent
\textbf{Contributions.} In this paper we show the equivalence of Russell and Tarski universes by a method due to \citet{Streicher},
for Martin Löf's Type Theory with dependent products, cumulative universes and judgmental $\beta\eta$ equality.
Our proof is modular and adapts to more complicated type systems. For instance, it adapts to a theory with universe polymorphism and internal universe levels \cite{polymorphism}.
More generally, we relate the Generalised Algebraic presentation of type theory to the implicit, and unannotated presentation
inherited from Pure Type Systems. In particular, our result justifies the practice of working in implicit
syntax, while legitimating the interpretation of implicit systems in usual models of type theory, such as in Categories with Families.
More precisely, our result proves that those two presentations have the same models.

The main argument is as follows: there is a natural forgetful map from the Tarski to the
Russell presentation. We show that any Russell judgement can be lifted in a unique way
to a Tarski judgement, proving the syntactic equivalence of the two theories.
This work is done in section \ref{sectionT2T2'}.
We rely on well-known properties of the syntax : no-confusion and injectivity of type formers, which are sometimes refered to as \textit{definitional inversion}.
We only need them at the level of the generalised algebraic theory, and they will propagate along our equivalence results.
Those properties seem to be strongly related to conversion \cite{meven}, and our equivalence result itself is to be
related to the elaboration in proof assistants.

Our proof is elementary and provides an algorithm to translate from one presentation to the other. Notably,
we only rely on the basic induction principles for the syntax, and not on stronger induction principles given
by normalisation. The definitional inversion properties are proven using the new domain-theoretic technique of \citet{definitional-inv}. Thus we believe that our proof could be adapted to various theories which do not enjoy this property.
In fact, the proof is purely syntactic and can be done internal to a weak metatheory such as primitive recursive
arithmetic, whereas normalisation relies on a stronger metatheory \cite{Coq19} such as the constructive set theory of \citet{aczel}.

In addition, we manually formalised and verified our equivalence result between Tarski and Russell universes in Rocq \cite{Formalisation}.
Throughout the paper, we provide links to the corresponding definitions and theorems
in the formalisation, indicated by the Rocq logo \includegraphics[height=1em]{icon-rocq-blue-eps-converted-to.pdf}.
Along with this, we did a complete auto-formalisation in Agda of the whole paper using Claude \cite{auto-formalisation}.\\

\noindent
\textbf{Structure of this paper.}
In Section \ref{GAT} we present the fully explicit Generalised Algebraic Theory and prove definitional inversion, before presenting the zoo of implicit presentations in Section \ref{implicit-presentations}.
Then, we prove that all those syntaxes are equivalent in the following sections: 

\begin{itemize}
    \item[] \textbf{\S\S~\ref{sectionT0T1}.} We relate explicit and implicit substitutions. 
    \item[] \textbf{\S\S~\ref{sectionT1T2}.} We relate Tarski and Russell presentations for cumulative universes.
    \item[] \textbf{\S\S~\ref{sectionT2T2'}.} We relate annotated and unannotated syntaxes. 
\end{itemize}

